% Memory hierarchy, map, caches, and the dmem mux. Requires tikz_styles.tex.
\begin{figure}[p]
\centering
\begin{tikzpicture}[node distance=4mm and 5mm]
  \node[ctl] (fe) {fetch};
  \node[ctl, below=of fe] (ms) {mem stage};
  \node[acc, below=of ms] (cont) {container};
  \node[acc, below=of cont] (sa) {STRACC};
  \node[acc, below=of sa] (gc) {GC};
  \node[rf, below=of gc] (spill) {RF spill};
  \node[ink, below=of spill] (fr) {frame / exc};
  \node[ex, below=of fr] (excore) {excore};

  \node[mem, right=14mm of ms, minimum width=26mm] (l1i) {L1I\\4-way, 8\,KB\\read-only};
  \node[mem, right=14mm of cont, minimum width=26mm] (l1d) {L1D\\4-way, 8\,KB\\write-back};
  \node[mem, right=12mm of l1d, minimum width=28mm] (l2) {L2 unified\\8-way, 128\,KB\\inclusive};
  \node[mem, right=8mm of l2, minimum width=24mm] (ram) {RAM\\16\,MB\\4-beat line};

  \draw[arr] (fe.east) -- (l1i.west);
  \draw[arr] (ms.east) -- ++(0.4,0) |- (l1d.west);
  \draw[arr] (cont.east) -- (l1d.west);
  \draw[arr] (sa.east) -- ++(0.35,0) |- (l1d.west);
  \draw[arr] (gc.east) -- ++(0.45,0) |- (l1d.west);
  \draw[arr] (spill.east) -- ++(0.55,0) |- (l1d.west);
  \draw[arr] (fr.east) -- ++(0.65,0) |- (l1d.west);
  \draw[arr] (l1i.east) -- ++(0.35,0) |- (l2.north west);
  \draw[arr] (l1d.east) -- (l2.west);
  \draw[arr] (l2.east) -- (ram.west);
  \draw[arr] (excore.east) -- ++(1.6,0) |- (l2.south west);

  \node[anno, anchor=north] at ($(excore.south)+(0.2,-0.15)$) {no L1D};
  \node[res, above=6mm of l1i] (codc) {CODC};
  \node[res, right=2mm of codc] (gic) {GIC};
  \node[res, right=2mm of gic, dashed] (ftb) {FTB\\not built};
  \draw[darr] (codc.south) -- ++(0,-0.25) -| (fe.north);
\end{tikzpicture}
\caption{Memory hierarchy behind an unchanged port.
A request is captured on the cycle \texttt{req} is high.
\texttt{ack} pulses for one cycle when the data is ready, any latency later, with \texttt{fault} beside it.
At most one ordinary request is outstanding per master.
\texttt{PYCORE\_CACHE\_EN=0} is a combinational bypass to the next level.
Line size is 64\,B everywhere.
L1 hit latency is 1 cycle; L2 hit latency is 8 cycles (\texttt{PYCORE\_L2\_HIT\_CYCLES}), a model parameter.
The crossbar adds \texttt{PYCORE\_CODE\_ADDR\_BASE} $= \mathtt{0x01000000}$ to instruction addresses so code and data do not alias in L2.
CODC and GIC sit beside the line caches.
The frame-top buffer (\texttt{PYCORE\_FTB\_FRAMES} $= 4$) is a named parameter only.}
\label{fig:hier}
\end{figure}

\begin{figure}[p]
\centering
\begin{tikzpicture}[x=1cm,y=1cm]
  % Data map, not to scale. Low addresses at the bottom.
  \begin{scope}[local bounding box=data]
    \fill[black!6] (0,0) rectangle (5.4,0.55);
    \fill[pyhi] (0,0.55) rectangle (5.4,1.05);
    \fill[pymem] (0,1.05) rectangle (5.4,4.55);
    \fill[pyink] (0,4.55) rectangle (5.4,5.15);
    \fill[pyctl] (0,5.15) rectangle (5.4,5.75);
    \fill[pyrf] (0,5.75) rectangle (5.4,6.55);
    \fill[pyacc] (0,6.55) rectangle (5.4,7.85);
    \draw[pyedge] (0,0) rectangle (5.4,7.85);
    \foreach \y in {0.55,1.05,4.55,5.15,5.75,6.55}
      \draw[pyedge] (0,\y) -- (5.4,\y);
    \node[font=\tiny, anchor=west] at (0.12,0.28) {reserved \texttt{0x0}};
    \node[font=\tiny, anchor=west] at (0.12,0.80) {boot record, 96\,B};
    \node[font=\tiny, anchor=west] at (0.12,2.80) {object heap, $\sim$15\,MB};
    \node[font=\tiny, anchor=west] at (0.12,4.85) {exc arena, 4\,KB};
    \node[font=\tiny, anchor=west] at (0.12,5.45) {frame stack, 32\,KB};
    \node[font=\tiny, anchor=west] at (0.12,6.15) {RF spill, 256\,KB};
    \node[font=\tiny, anchor=west] at (0.12,7.20) {GC metadata, 512\,KB};
    \node[anno, anchor=south] at (2.7,8.0) {data, not to scale};
  \end{scope}
  \begin{scope}[xshift=7.2cm, local bounding box=code]
    \fill[pymem] (0,0) rectangle (5.6,2.3);
    \fill[pyhi] (0,2.3) rectangle (5.6,5.4);
    \draw[pyedge] (0,0) rectangle (5.6,5.4);
    \draw[pyedge] (0,2.3) -- (5.6,2.3);
    \node[font=\tiny, align=center] at (2.8,1.15)
      {code ROM\\slots \texttt{0x0000}--\texttt{0x1FFF}\\64\,KB, read-only};
    \node[font=\tiny, align=center] at (2.8,3.85)
      {code RAM\\slots \texttt{0x2000}--\texttt{0x21FFF}\\1\,MB, writable};
    \node[anno, anchor=south] at (2.8,5.55) {code, after $+$\texttt{0x01000000}};
  \end{scope}
  \node[font=\tiny\sffamily, align=left, anchor=north west] at (0,-0.35)
    {Heap base \texttt{0x440}. Heap limit \texttt{0xF00000}.\\
     Exc \texttt{0xF00000}. Frames \texttt{0xF01000}. Spill \texttt{0xF40000}.\\
     GC meta \texttt{0xF80000}--\texttt{0xFFFFFF}. Data limit $=$ code base.};
\end{tikzpicture}
\caption{Address map.
The heap is a bump pointer from \texttt{PYCORE\_HEAP\_BASE} with no free list until the collector is enabled.
The static boot image occupies about 380\,KB of that heap.
Code is a separate slot space: the architectural PC is a slot index, and fetch sends \texttt{pc}$\ll 3$.
\texttt{pycore\_code\_mem} is compiled and not instantiated; the live path is \texttt{pycore\_mem\_hier} into \texttt{pycore\_ram}.}
\label{fig:map}
\end{figure}

\begin{figure}[htbp]
\centering
\begin{tikzpicture}[node distance=3.5mm and 4mm]
  \node[bus] (req) {addr};
  \node[mem, right=of req] (idx) {set index};
  \node[mem, right=of idx] (tag) {tag compare\\$N$ ways};
  \node[rf, right=of tag] (lru) {LRU};
  \node[mem, below=of tag] (hit) {hit data};
  \node[acc, left=of hit] (miss) {miss};
  \node[mem, below=of miss] (fill) {line fill\\64\,B};
  \node[trap, right=of fill] (wbv) {dirty victim\\writeback};
  \node[ctl, right=of hit] (ack) {ack\\1 cycle};

  \draw[arr] (req) -- (idx);
  \draw[arr] (idx) -- (tag);
  \draw[arr] (tag) -- (lru);
  \draw[arr] (tag) -- (hit);
  \draw[arr] (tag) -- (miss);
  \draw[arr] (miss) -- (fill);
  \draw[arr] (fill) -- (wbv);
  \draw[arr] (hit) -- (ack);
  \draw[arr] (fill.east) -- ++(0.3,0) |- (ack.south);
\end{tikzpicture}
\caption{One \texttt{pycore\_cache}, three instantiations (L1I, L1D, L2).
L1I is read-only and write-invalidate / no-allocate: a store drops a hit line and forwards the write, and does not allocate.
L1D and L2 are write-back and write-allocate.
An all-zero full-line write on L1D bypasses allocation; STRACC uses full-line installs so a new result line is not filled only to be overwritten.
\texttt{rdata\_line} returns the whole responding line with \texttt{ack}, which is how the fetch buffer fills.}
\label{fig:cache}
\end{figure}

\begin{figure}[htbp]
\centering
\begin{tikzpicture}[x=1cm,y=1cm]
  \node[acc, minimum width=32mm] (nb) at (0,1.6) {L1D non-blocking\\up to 4 fills};
  \node[mem, minimum width=32mm] (hold) at (4.6,1.6) {ordinary miss\\one-entry hold};
  \node[mem, minimum width=28mm] (inst) at (0,0) {install oldest};
  \node[rf, minimum width=28mm] (wbb) at (4.6,0) {writeback\\one line};
  \node[mem, minimum width=36mm] (pipe) at (9.0,0.8) {L2 pipe port\\in order, 4 beats};
  \draw[arr] (nb) -- (hold);
  \draw[arr] (nb) -- (inst);
  \draw[arr] (inst) -- (wbb);
  \draw[arr] (hold) -- (wbb);
  \draw[arr] (wbb) -- (pipe);
  \draw[arr] (hold.east) -- ++(0.35,0) |- (pipe.north);
\end{tikzpicture}
\caption{The two opt-in ports.
Nothing changes for a master that ignores them: an ordinary request takes the same number of cycles as before.
The collector is the first user.
While marking, it prefetches the line of each object it pushes and the next lines of the range it is scanning.
A prefetch fills L1D and does not produce an answer.}
\label{fig:nb}
\end{figure}

\begin{figure}[htbp]
\centering
\begin{tikzpicture}[node distance=2.4mm]
  \node[ex, minimum width=52mm] (a) {recoverable trap detected};
  \node[mem, minimum width=52mm, below=of a] (b) {L1D writeback $+$ invalidate};
  \node[ex, minimum width=52mm, below=of b] (c) {mailbox accepts \texttt{trap\_req}};
  \node[ex, minimum width=52mm, below=of c] (d) {grant flips to excore};
  \node[ctl, minimum width=52mm, below=of d] (e) {firmware completes the effect};
  \node[mem, minimum width=52mm, below=of e] (f) {L1D invalidate};
  \node[ctl, minimum width=52mm, below=of f] (g) {\texttt{trap\_res} to PyCore};
  \node[mem, minimum width=52mm, below=of g] (h) {flush CODC and GIC};
  \node[ctl, minimum width=52mm, below=of h] (i) {grant returns to PyCore};
  \foreach \p/\q in {a/b,b/c,c/d,d/e,e/f,f/g,g/h,h/i}
    \draw[arr] (\p) -- (\q);
\end{tikzpicture}
\caption{The only sharing event.
There is no cache coherence protocol.
PyCore issues no data request while excore owns memory, and excore polls the mailbox while PyCore owns it.
L1I is not flushed on the handoff.
L2 is left in place; excore's writes land there and become visible to PyCore after the L1D invalidate.
A code-RAM blit or patch invalidates L1I and flushes CODC.
\texttt{STORE\_NAME} / \texttt{STORE\_GLOBAL} and a change of \texttt{globals\_base\_r} flush the GIC.
A return that restores the same globals base does not.}
\label{fig:handoff}
\end{figure}

\begin{figure}[htbp]
\centering
\begin{tikzpicture}[node distance=3mm and 10mm]
  \node[mem, minimum width=48mm, align=left, font=\tiny\sffamily] (codc)
    {\textbf{CODC} \quad 4 entries, 2 ways\\
     key: code-object base\\
     set: \texttt{addr}$\gg 6$\\
     payload 576 bits:\\
     entry slot, consts, names,\\
     metadata, defaults};
  \node[mem, minimum width=48mm, align=left, font=\tiny\sffamily, right=of codc] (gic)
    {\textbf{GIC} \quad 16 entries, 2 ways\\
     key: \{\,code, namei\,\}\\
     set: \texttt{namei[2:0]}\\
     payload: one 132-bit entry\\
     fill only on a hit resolve\\
     a miss is never cached};
  \draw[darr] ($(codc.south)+(0,-0.15)$) -- ++(0,-0.01);
\end{tikzpicture}
\caption{Result caches.
Both answer in the same cycle as the lookup.
Fill and LRU update on the following clock edge.
\texttt{CACHE\_EN=0} forces a miss.
CODC is flushed on \texttt{MAKE\_FUNCTION}, code release, \texttt{trap\_res}, and reset.
The GIC is flushed as a whole on purpose: consulting a dict version to do better would itself be a data-memory access.}
\label{fig:resultcache}
\end{figure}

\begin{figure}[htbp]
\centering
\begin{tikzpicture}[node distance=3.5mm and 4mm]
  \node[mem] (l1i) {L1I line\\64\,B, 8 slots};
  \node[ctl, right=of l1i] (buf) {fetch line buffer};
  \node[ctl, right=of buf] (fold) {skip CACHE\\fold EXTENDED\_ARG};
  \node[ctl, right=of fold] (pc) {architectural PC\\one slot};
  \draw[arr] (l1i) -- (buf);
  \draw[arr] (buf) -- (fold);
  \draw[arr] (fold) -- (pc);
  \node[anno, align=left, anchor=north west] at ($(l1i.south west)+(0,-0.35)$)
    {Later slots in the same line, including the skips, do not issue another imem request.\\
     L1I hit and miss counters therefore count fills behind this buffer, not wordcode slots.\\
     A code-RAM write drops the buffer so a later fetch cannot observe a stale slot.};
\end{tikzpicture}
\caption{Fetch.
The image keeps CPython's \texttt{CACHE} and \texttt{EXTENDED\_ARG} bytes.
The slot index equals the CPython code-unit index, so branch arguments are not rewritten.
The live code image is ROM plus RAM in the unified RAM, selected by the slot address.}
\label{fig:fetch}
\end{figure}

\begin{figure}[htbp]
\centering
\begin{tikzpicture}[node distance=1.6mm]
  \node[rf, minimum width=62mm] (s1) {1 \quad RF spill / fill};
  \node[ink, minimum width=62mm, below=of s1] (s2) {2 \quad frame push / pop};
  \node[acc, minimum width=62mm, below=of s2] (s3) {3 \quad container, boot, code-object fields};
  \node[acc, minimum width=62mm, below=of s3] (s4) {4 \quad STRACC};
  \node[ink, minimum width=62mm, below=of s4] (s5) {5 \quad exception stack};
  \node[acc, minimum width=62mm, below=of s5] (s6) {6 \quad GC engine};
  \node[acc, minimum width=62mm, below=of s6] (s7) {7 \quad GC run allocator};
  \node[ctl, minimum width=62mm, below=of s7] (s8) {8 \quad mem stage (pointer ops)};
  \node[mem, right=8mm of s4] (port) {one\\dmem\\port};
  \foreach \n in {s1,s2,s3,s4,s5,s6,s7,s8}
    \draw[arr] (\n.east) -- (port.west);
\end{tikzpicture}
\caption{Who drives the data port.
The mux is a priority encode, but the owners are mutually exclusive by core state.
PyCore never arbitrates two live masters in one cycle.
The port is 128-bit data, 16-byte aligned, with a 16-bit write strobe.
A misaligned or out-of-range access raises \texttt{ADDR\_ALIGN} or \texttt{MEM\_FAULT}.
GC also has the non-blocking prefetch port into L1D, which is not this mux.}
\label{fig:dmux}
\end{figure}
